\documentclass[11pt]{article}
\usepackage[margin=25mm]{geometry}
\usepackage{fontspec}
\setmainfont{Latin Modern Roman}
\usepackage{amsmath,amssymb,mathtools}
\usepackage{unicode-math}
\setmathfont{Latin Modern Math}
\usepackage{xurl,hyperref}
\hypersetup{colorlinks=true,urlcolor=blue}
\setcounter{secnumdepth}{0}
\setlength{\parindent}{0pt}
\setlength{\parskip}{5pt}
\setlength{\emergencystretch}{3em}
\newcommand{\sref}[2]{\hyperlink{source:#1:#2}{[S#2]}}
\newcommand{\cataloguescope}[1]{\paragraph{Recorded target.}#1}
\newcommand{\R}{\mathbb R}
\newcommand{\C}{\mathbb C}
\newcommand{\N}{\mathbb N}
\newcommand{\Z}{\mathbb Z}
\newcommand{\Q}{\mathbb Q}
\newcommand{\F}{\mathbb F}
\newcommand{\eps}{\varepsilon}
\begin{document}
% UnsolvedMath ID 20002537; source catalogue code AIM-PHYSICS-0045
\setcounter{section}{20002536}
\section{Syzygy-word parity and the angular-momentum holonomy obstruction}\label{20002537}
{\small UnsolvedMath ID 20002537 · Mathematical Physics}
\subsection{Definitions and mathematical statement}

\paragraph{Shared notation.}$\mathbb N=\{1,2,\ldots\}$; the standard number fields have their usual meanings. All additional conventions are stated locally.


\paragraph{Definitions and notation.}In the Newtonian $n$-body problem, positions $q_i\in\mathbb R^d$ and masses $m_i>0$ obey $m_i\ddot q_i=\nabla_{q_i}U$, where $U(q)=\sum_{i<j}m_im_j/|q_i-q_j|$ and the gravitational constant is normalized to one. Here $T=\tfrac12\sum_i m_i|\dot q_i|^2$, $E=T-U$, and a collision has $q_i=q_j$. The center of mass can be fixed at zero; moment of inertia is $I=\sum_i m_i|q_i|^2$. For a fixed period $P$, the action is $\mathcal A(q)=\int_0^P(T+U)\,dt$. A collision-free minimizer in an admissible loop class satisfies the Newtonian equations when variations are permitted. The period, constraints and closure of the class must be specified: a minimizing sequence can escape or approach collisions. For the planar three-body problem, a shape is a configuration modulo translation, rotation and positive scaling. Its normalized shape space is a sphere, with collisions marked and collinear configurations on an equator. A syzygy is a collinear configuration; its symbol records which body is in the middle. 

\paragraph{Catalogue wording retained for traceability.}
Can every symbol sequence, obtained by recording successive collinear configurations with the middle body as symbol, be realized in the planar three-body problem? Determine whether the realizations minimize action and whether they have zero angular momentum.
\paragraph{Scope and source qualification.}The source does not fix finite versus infinite or periodic symbol sequences, repetitions, endpoints or collision treatment. Such choices must be supplied in a theorem and identified as a clarified scope, not silently asserted for all conceivable sequences. The catalogue title/status was paired with a machine-generated research attempt; that attempt is not peer-reviewed evidence of settlement. The statement covers the original source questions rather than only the title’s subcase.
\subsection{Short English statement}
Determine which three-body syzygy symbol sequences are realizable, minimizing, and compatible with zero angular momentum.
\subsection{Research attempt}
\paragraph{Model, date and investigated scope.}
OpenAI GPT-6 Astra (configured model \texttt{gpt-6-astra}), ultra reasoning; generated 11 October 2026. Proved even crossing parity for transverse periodic syzygies and derived the zero-angular-momentum reconstruction angle. General symbolic realization, minimizing behavior and zero-angular-momentum existence remain unresolved.

\paragraph{Sources and literature check.}
Read the supplied complete definition, catalogue wording, scope qualifications and source list for this canonical ID. The original-source citations are inherited as supplied references; no fresh primary-source or novelty search was performed for this attempt. No current open-status claim is inferred from this calculation.

\paragraph{A parity obstruction for one clarified sequence class.}For a collision-free planar three-body periodic motion, let \(A(t)=\det(q_2-q_1,q_3-q_1)\) be twice the oriented area. Suppose all syzygies in a period are isolated transverse zeros of \(A\), and the chosen period endpoints are not syzygies. At each such zero the sign changes. Since \(A(P)=A(0)\), the number of crossings in the period must be even. Thus a periodic word of odd length cannot be realized under these transversality and endpoint conventions. Tangential contacts, persistent collinearity or a different coding rule invalidate this particular count and are not silently excluded from the source's broad question.
\paragraph{Zero angular momentum does not guarantee a closed lift.}Write \(q_i=r(t)e^{i\theta(t)}s_i(t)\), where \(r>0\) is scale and \(s_i\) is a local normalized shape representative, with \(I_s=\sum m_i|s_i|^2>0\). Direct differentiation gives
\[C=r^2\left(I_s\dot\theta+\operatorname{Im}\sum_i m_i\overline{s_i}\dot s_i\right).
\]
At \(C=0\), integrating yields \(\theta(P)-\theta(0)=-\int_0^P I_s^{-1}\operatorname{Im}\sum_i m_i\overline{s_i}\dot s_i\,dt\). Even if the chosen shape representative is periodic, closure of the physical configuration requires this angle to lie in \(2\pi\mathbb Z\). The syzygy word alone does not determine the integral.
\paragraph{Gap.}These are necessary compatibility tests, not a realization theorem. The first missing step is a collision-free trajectory with prescribed coding and controlled reconstruction angle. Action minimization, zero-angular-momentum realization, finite versus infinite sequences and repeated/tangent symbols all require additional analysis. The broad source ambiguity is retained rather than declaring every odd symbolic sequence impossible in every convention.

\paragraph{Verification and adversarial audit.}
Checked oriented-area sign changes, periodic endpoint assumption, tangencies, rotation derivative and 2 pi closure requirement. This was an LLM adversarial audit, not independent human verification; no theorem prover was used. Prompts A then B applied to this ID: Checked oriented-area sign changes, periodic endpoint assumption, tangencies, rotation derivative and 2 pi closure requirement. Scope remains limited: Construct collision-free coded trajectories while controlling the reconstruction holonomy and the intended finite/infinite coding convention. Final pass clarified the separate positive scale factor in angular-momentum reconstruction.

Coordinator checked the transverse syzygy parity condition and zero-angular-momentum reconstruction-angle identity. This additional check was model-assisted, not independent human verification.

\paragraph{Tools, process and contributions.}
Python 3 and local shell used to read supplied files and serialize this manuscript; mathematics derived symbolically in the text, with no numerical experiment or formal proof assistant. Contribution prompt A was followed by prompt B. GPT-6 Astra with ultra reasoning drafted the research and performed its stated self-audit; any additional coordinator contribution is identified in the verification and audit. A GPT-6-based Codex coordinator assembled the manuscript and checked repository integration; its exact variant and reasoning configuration were not exposed. No human mathematical contribution, independent human proof check or formal proof certification is supplied. Complete derivations and the consulted source references are included here; private reasoning transcripts are not published.

\paragraph{Final status and first remaining gap.}
\textbf{UNRESOLVED}. Construct collision-free coded trajectories while controlling the reconstruction holonomy and the intended finite/infinite coding convention.

\section{Completion Estimate}
\noindent\textbf{COMPLETION: 0\%}
The calculation is an elementary benchmark or conditional reduction; it establishes no claimed new portion of the complete source target.

\subsection{Sources}
\begin{itemize}
\item[\hypertarget{source:20002537:1}{S1}] ULAM.AI and the UnsolvedMath Contributors, UnsolvedMath, record 20002537 (AIM-PHYSICS-0045). Catalogue curation under CC BY 4.0; definitions and English exposition added locally.
\url{https://huggingface.co/datasets/ulamai/UnsolvedMath}.
\item[\hypertarget{source:20002537:2}{S2}] American Institute of Mathematics, Variational Methods in Celestial Mechanics; notes by Kuo-Chang Chen, version 22 June 2003, source item 8.
\url{https://aimath.org/WWN/varcelest/varcelest.pdf}.
\end{itemize}
\label{end:20002537}
\end{document}
